\subsection{Independent proof review of four-window spectral sharpness}\label{endpoint-four_window_sharpness_review--independent-proof-review-of-four-window-spectral-sharpness}

Date: 13 September 2026. Accepted source: work/tau\_four\_window\_spectral\_sharpness\_20260913.md, 18872 bytes, SHA-256 11c26c3bcfce135e1882190af9dc3a8f1e4e6487fd282cc252fa441fdf90651c. The complete FS1--FS35 proof was read. The projection and strictness arguments, the explicit coefficient-metric family, its control and difference spectra, and its stated arithmetic scope are accepted without correction.

The source relies on the finite spaces of work/tau\_arithmetic\_determinant\_transport\_20260913.tex, 17706 bytes, SHA-256 24cf7ce5be6a43cc656641b516185df10b2e4d7824600ece657120d1c549deae. Those exact spaces, original source maps, and determinant transport were independently verified in work/tau\_arithmetic\_determinant\_transport\_independent\_review\_20260913.md, SHA-256 ce6b9fbdad43e5233b76756f5cbbbadf6aca0a33db62437f7caf0b90321683bb.

No closed AW, WS, or AT source or PDF was changed. This review is an additional complete proof body. It includes an exact refinement of the supplied family: its fixed-metric optimum over the prescribed control spectrum is computed explicitly below. No old test suite, Lean process, numerical moment calculation, or remote operation was used.

\subsubsection{1. Literal flags, dimensions, and two exact individual spectra}\label{endpoint-four_window_sharpness_review--literal-flags-dimensions-and-two-exact-individual-spectra}

Fix the original monic polynomial \(\chi\in\mathbb C[S]\) of degree \(q\ge1\). Set \(n=2q+1\), \(\mathcal H=\mathcal P_{2q}\), and retain an arbitrary positive coefficient form \(M\), with inner product \(x^*My\). All projections and complements in this review use that same form. Let \(P_N\) project onto \(\mathcal P_N\), and \(Q_N\) onto \(\chi\mathcal P_{N-q}\), with \(Q_{q-1}=0\). Define

\[
U=Q_{2q-1}+Q_{2q}-Q_q,\quad
W=P_{2q-1}-P_{q-1}+I-P_q,\quad D=U-W.
\tag{SR1}
\]

For nested subspaces \(X\subset Y\), their orthogonal projections have products \(P_XP_Y=P_YP_X=P_X\). Indeed the orthogonal decomposition \(X\oplus(Y\cap X^\perp)\oplus Y^\perp\) gives their diagonal actions. Thus each of the individual degree and relation flags is a commuting projection flag.

Set \(B_0=\mathbb C\chi\), \(B_1=\chi\mathcal P_{q-1}\), \(B_2=\chi\mathcal P_q\), of dimensions \(1,q,q+1\). These dimensions use injectivity of multiplication by the original nonzero polynomial. The relation flag decomposes \(\mathcal H\) orthogonally into \(B_0\), \(B_1\cap B_0^\perp\), \(B_2\cap B_1^\perp\), and \(B_2^\perp\). Substitution in SR1 shows that \(U\) has respective values \(1,2,1,0\). The degree flag similarly decomposes into \(\mathcal P_{q-1}\), \(\mathcal P_q\cap\mathcal P_{q-1}^\perp\), \(\mathcal P_{2q-1}\cap\mathcal P_q^\perp\), and \(\mathcal P_{2q-1}^\perp\); \(W\) has respective values \(0,1,2,1\). Consequently both operators have the precise spectrum

\[
0^{[q]},\quad 1^{[2]},\quad 2^{[q-1]},
\qquad \operatorname{Tr}U=\operatorname{Tr}W=2q,\quad
\operatorname{Tr}D=0.
\tag{SR2}
\]

For \(q=1\), the multiplicity of two is zero; \(B_0=B_1\), and \(U=Q_2\). These are sums of projections with the indicated weights; for \(q>1\) neither is an idempotent.

\subsubsection{2. Trace rearrangement with all equality conditions needed later}\label{endpoint-four_window_sharpness_review--trace-rearrangement-with-all-equality-conditions-needed-later}

Let \(A\) be self-adjoint with ordered eigenvalues \(a_1\le\cdots\le a_n\). For a rank-\(r\) orthogonal projection \(E\), choose an orthonormal eigenbasis \(v_i\), and write \(e_i=\langle v_i,Ev_i\rangle=\|Ev_i\|^2\). Then \(0\le e_i\le1\), \(\sum_i e_i=r\), and \(\operatorname{Tr}(AE)=\sum_i a_ie_i\).

The largest possible sum of these weights is the sum of the largest \(r\) eigenvalues. To see this directly, choose a threshold \(a_{n-r+1}\). The mass on indices strictly below the threshold has nonpositive contribution after subtracting that threshold times \(r\), and the mass above it contributes at most the sum of their full coefficients. Assigning weight one at the top \(r\) indices reaches the resulting bound. The analogous argument at the bottom gives

\[
\sum_{i=1}^{r}a_i\le\operatorname{Tr}(AE)
\le\sum_{i=n-r+1}^{n}a_i.
\tag{SR3}
\]

Equality on the upper side forces \(E\) to contain every eigenspace with eigenvalue strictly above the cutoff and be orthogonal to every eigenspace strictly below it. A diagonal value one gives \(\|(I-E)v_i\|^2=0\); a value zero gives \(Ev_i=0\). The only free subspace is inside the eigenspace at the cutoff. The lower side has the reversed conditions. Thus this proof retains arbitrary repeated eigenvalues.

For a self-adjoint \(C\) with ordered eigenvalues \(c_1\le\cdots\le c_n\), decompose \(U\) as its positive spectral projection plus its level-two projection, of ranks \(q+1\) and \(q-1\). Applying SR3 and doing the same for \(W\) yields

\[
\begin{aligned}
\operatorname{Tr}(UC)&\le c_{q+1}+c_{q+2}
+2\sum_{i=q+3}^{2q+1}c_i,\\
\operatorname{Tr}(WC)&\ge c_q+c_{q+1}
+2\sum_{i=1}^{q-1}c_i.
\end{aligned}
\tag{SR4}
\]

Subtracting, and also reversing the two extrema, proves

\[
|\operatorname{Tr}(DC)|\le B_q(c),\qquad
B_q(c)=c_{q+2}-c_q+
2\sum_{j=1}^{q-1}(c_{2q+2-j}-c_j).
\tag{SR5}
\]

For \(q=1\), this is \(c_3-c_1\). If only the two individual spectra in SR2 are prescribed, equality is attained by diagonal operators whose ordered entries are \(u=(0^{[q]},1,1,2^{[q-1]})\) and its reversal, in an eigenbasis of \(C\). Their difference has ordered eigenvalues

\[
\lambda^*=(-2^{[q-1]},-1,0,1,2^{[q-1]}).
\tag{SR6}
\]

Direct multiplication of the diagonal entries gives precisely SR5 with positive sign; exchanging the two operators gives the negative sign. This proves the claimed sharp value for that class of two independent spectral orbits.

\subsubsection{3. The polynomial intersections exclude every nontrivial rank equality}\label{endpoint-four_window_sharpness_review--the-polynomial-intersections-exclude-every-nontrivial-rank-equality}

The two needed positions of the literal polynomial subspaces follow solely from degrees:

\[
\mathcal P_{q-1}\cap B_2=0,\qquad
\mathcal P_q+B_2=\mathcal H.
\tag{SR7}
\]

A nonzero multiple of \(\chi\) has degree at least \(q\), proving the first. In the second, \(B_2\cap\mathcal P_q=\mathbb C\chi\), because multiplication by a nonconstant polynomial has degree exceeding \(q\). Its sum therefore has dimension \((q+1)+(q+1)-1=2q+1\), proving equality with the whole space.

Let \(U_j=\ker(U-jI)\), \(W_j=\ker(W-jI)\). The eigenspaces computed in Section 1 and SR7 imply

\[
U_2\cap W_0=0,\qquad W_2\cap U_0=0.
\tag{SR8}
\]

The first uses \(U_2\subset B_2\) and \(W_0=\mathcal P_{q-1}\). For the second, a vector in \(W_2\cap U_0\) is orthogonal to both \(\mathcal P_q\) and \(B_2\), hence to their whole sum.

Let \(a_r\) be the sum of the largest \(r\) entries of SR6, for \(1\le r<n\). Applying SR3 separately to \(U\) and \(W\) gives, for every rank-\(r\) orthogonal projection \(E\),

\[
\operatorname{Tr}(DE)\le
\sum_{i=n-r+1}^{n}u_i-\sum_{i=1}^{r}u_i=a_r.
\tag{SR9}
\]

Equality in this subtraction requires equality in both component bounds; their two deficits are nonnegative. The exact equality conditions from SR3 now give the following exhaustive cases.

For \(1\le r\le q-1\), the maximum for \(U\) is \(2r\) and the minimum for \(W\) is zero. Equality therefore requires \(\operatorname{im}E\subset U_2\cap W_0\), excluded by SR8.

For \(r=q\) and \(q\ge2\), minimization for \(W\) forces \(\operatorname{im}E=W_0\), because that kernel already has dimension \(q\). Maximization for \(U\) forces the image to contain its entire level-two space of dimension \(q-1>0\). This contradicts \(U_2\cap W_0=0\). At \(q=r=1\), the maximum for \(U\) instead says \(\operatorname{im}E\subset B_2\), while the minimum for \(W\) says \(\operatorname{im}E=\mathcal P_0\); SR7 excludes this case as well.

For \(r=q+1\), the maximum for \(U\) forces the image to equal its entire positive spectral space \(B_2\): its dimension is exactly \(q+1\). The minimum for \(W\) forces the image to contain \(W_0=\mathcal P_{q-1}\). Since \(q\ge1\), that subspace is nonzero, contradicting SR7.

For \(q+2\le r\le2q\), put \(F=I-E\), of rank \(s=n-r\) with \(1\le s\le q-1\). Since \(\operatorname{Tr}D=0\) and the spectrum SR6 is invariant under negation, \(a_r=a_s=2s\). Equality in SR9 would force \(\operatorname{Tr}(DF)=-2s\). The lower bound for \(U\) and upper bound for \(W\) then require \(\operatorname{im}F\subset U_0\cap W_2\), again excluded by SR8.

These arguments cover every rank including \(q=1\), and use the actual repeated eigenspaces of \(U,W\); they do not suppose that \(C\) has simple spectrum.

For a fixed positive \(M\), the rank-\(r\) orthogonal projections form a compact set. Conjugation by the explicit isometry \(M^{1/2}\) identifies it with the Euclidean self-adjoint idempotents of trace \(r\), a closed bounded set in finite matrix dimension. Hence the continuous functional \(\operatorname{Tr}(DE)\) attains its maximum. Since equality in SR9 is impossible, if \(d_1\le\cdots\le d_n\) are the actual eigenvalues of \(D\), SR3 applied to \(D\) gives

\[
\sum_{i=n-r+1}^{n}d_i<a_r\quad(1\le r<n),
\qquad \sum_i d_i=0.
\tag{SR10}
\]

\subsubsection{4. Strictness for every nonscalar control, with repeated eigenvalues retained}\label{endpoint-four_window_sharpness_review--strictness-for-every-nonscalar-control-with-repeated-eigenvalues-retained}

Choose one orthonormal eigenbasis of \(C\), in any order inside a repeated eigenspace, and let \(E_r\) project onto its last \(r\) vectors. The exact expansion is

\[
C=c_1I+\sum_{r=1}^{n-1}
(c_{n-r+1}-c_{n-r})E_r.
\tag{SR11}
\]

All coefficients are nonnegative. The scalar term has zero trace against \(D\). Applying SR10 term by term therefore yields

\[
\operatorname{Tr}(DC)
<\sum_{r=1}^{n-1}(c_{n-r+1}-c_{n-r})a_r=B_q(c)
\tag{SR12}
\]

whenever \(C\) is nonscalar: at least one spectral gap is then strictly positive, and its corresponding inequality in SR10 is strict. The last equality follows by telescoping with the ordered zero-sum vector SR6, or by multiplying its entries with \(c_i\). Applying the same argument to \(-C\) proves the negative bound; its ordered eigenvalues are \(-c_n,\ldots,-c_1\), and the symmetry of SR6 gives the same \(B_q(c)\). Hence

\[
|\operatorname{Tr}(DC)|<B_q(c)
\quad\text{for every nonscalar \(C\) and every fixed \(M>0\)}.
\tag{SR13}
\]

For scalar \(C\), both sides are zero. Every \(a_r\) is strictly positive, so SR11 also proves \(B_q(c)>0\) for every nonscalar \(C\). This completes the repeated-spectrum part of FS19.

The finite strictness has an exact quantitative record: define \[
g_r(M)=a_r-\sum_{i=n-r+1}^{n}d_i>0.
\tag{SR14}
\] Before making SR12 strict, its positive-side deficit satisfies \[
B_q(c)-\operatorname{Tr}(DC)
\ge\sum_{r=1}^{n-1}(c_{n-r+1}-c_{n-r})g_r(M)>0.
\tag{SR15}
\] For the negative side use the expansion of \(-C\) and its own ordered gaps against the same \(g_r(M)\). This retains every spectral gap and proves exactly where strictness comes from. It supplies no uniform lower bound for those positive deficits as \(M\) varies.

\subsubsection{5. The complete metric family on the unchanged polynomial flag}\label{endpoint-four_window_sharpness_review--the-complete-metric-family-on-the-unchanged-polynomial-flag}

Retain every coefficient of the original \(\chi\). The ordered polynomials

\[
b_i=S^i\ (0\le i<q),\quad b_q=\chi,\quad
b_{q+j}=\chi S^j\ (1\le j\le q)
\tag{SR16}
\]

have increasing degrees and leading coefficient one. Their coordinate matrix relative to the monomial frame is triangular with determinant one, and each initial segment spans the original degree space.

Fix \(\mu>0,\epsilon>0\). Define \(T_\epsilon:\mathcal H\to\mathbb C^n\) by \[
T_\epsilon b_i=e_i\ (0\le i\le q),\qquad
T_\epsilon b_{q+j}=e_{j-1}+\epsilon e_{q+j}\ (1\le j\le q).
\tag{SR17}
\] It has determinant \(\epsilon^q\), and its inverse sends \(e_i\) to \(b_i\) for \(i\le q\), and \(e_{q+j}\) to \(\epsilon^{-1}(b_{q+j}-b_{j-1})\). Substituting each basis vector into these two formulas verifies both inverse compositions.

Set \(M_\epsilon=\mu T_\epsilon^*T_\epsilon\) as a form on the original polynomial coefficients. The explicit isometry to the Euclidean target is \(\sqrt\mu\,T_\epsilon\). For \(P=\sum_i x_ib_i\), its full norm is

\[
\|P\|_{M_\epsilon}^2
=\mu\left(
|x_q|^2+\sum_{j=1}^q|x_{j-1}+x_{q+j}|^2
+\epsilon^2\sum_{j=1}^q|x_{q+j}|^2\right).
\tag{SR18}
\]

The cross term in each two-coordinate block is retained. Each such block is \(\mu\begin{pmatrix}1&1\\1&1+\epsilon^2\end{pmatrix}\), of determinant \(\mu^2\epsilon^2\), and the remaining \(b_q\) block is \(\mu\). Therefore

\[
\det M_\epsilon=\mu^{2q+1}\epsilon^{2q},\quad
\|1\|^2=\|\chi\|^2=\mu,\quad
\|\chi S^j\|^2=\mu(1+\epsilon^2),\quad
\langle S^{j-1},\chi S^j\rangle=\mu.
\tag{SR19}
\]

The determinant is the same in the original monomial frame, by the unit determinant of SR16. All unlisted cross terms between different paired blocks or the \(b_q\) block are zero.

The map SR17 transports the degree flag to the standard coordinate flag. Thus \(W\) has values zero on \(e_0,\ldots,e_{q-1}\), one on \(e_q,e_{2q}\), and two on \(e_{q+1},\ldots,e_{2q-1}\). Define the orthonormal vectors \[
v_j=\frac{e_{j-1}+\epsilon e_{q+j}}{\sqrt{1+\epsilon^2}}
\quad(1\le j\le q).
\tag{SR20}
\] They are orthogonal to one another and to \(e_q\). The transformed relation spaces are spanned by \(e_q\), then \(e_q,v_1,\ldots,v_{q-1}\), then \(e_q,v_1,\ldots,v_q\). Their projections therefore give \[
U=|e_q\rangle\langle e_q|
+2\sum_{j=1}^{q-1}|v_j\rangle\langle v_j|
+|v_q\rangle\langle v_q|.
\tag{SR21}
\] The unit vectors in SR20 describe the coordinate isometry; the original masses and norms remain SR18--SR19.

\subsubsection{6. Prescribed spectrum, exact trace, and full difference spectrum}\label{endpoint-four_window_sharpness_review--prescribed-spectrum-exact-trace-and-full-difference-spectrum}

For prescribed ordered real numbers \(c_1,\ldots,c_n\), set \[
a_j=c_{2q+2-j},\ b_j'=c_j\quad(1\le j<q),\qquad
a_q=c_{q+2},\ b_q'=c_q.
\tag{SR22}
\] Give \(e_{j-1}\) eigenvalue \(a_j\), \(e_{q+j}\) eigenvalue \(b_j'\), and \(e_q\) eigenvalue \(c_{q+1}\). This diagonal \(C_0\) uses every prescribed eigenvalue exactly once. Its pullback \(C_\epsilon=T_\epsilon^{-1}C_0T_\epsilon\) has, on each ordered pair \((b_{j-1},b_{q+j})\), matrix \[
\begin{pmatrix}a_j&a_j-b_j'\\0&b_j'\end{pmatrix}.
\tag{SR23}
\] In particular its entries are independent of \(\epsilon\). Multiplication by the original Gram block gives \[
\mu\begin{pmatrix}1&1\\1&1+\epsilon^2\end{pmatrix}
\begin{pmatrix}a_j&a_j-b_j'\\0&b_j'\end{pmatrix}
=\mu\begin{pmatrix}a_j&a_j\\a_j&a_j+\epsilon^2b_j'\end{pmatrix},
\tag{SR24}
\] which is Hermitian because \(a_j,b_j'\) are real. This proves self-adjointness in every \(M_\epsilon\), retaining all masses and cross terms.

In the isometric coordinates, each paired block of \(D\) is \[
D_j=\frac{w_j}{1+\epsilon^2}
\begin{pmatrix}1&\epsilon\\\epsilon&-1\end{pmatrix},\qquad
w_j=2\ (j<q),\quad w_q=1.
\tag{SR25}
\] The common coordinate \(e_q\) has value zero. The diagonal \(C_0\) therefore contributes \(w_j(a_j-b_j')/(1+\epsilon^2)\) to the trace of each block. Summing proves \[
\operatorname{Tr}(DC_\epsilon)=\frac{B_q(c)}{1+\epsilon^2}.
\tag{SR26}
\]

The unweighted matrix in SR25 has trace zero and determinant \(-1/(1+\epsilon^2)\). Its two eigenvalues are \(\pm1/\sqrt{1+\epsilon^2}\). Including the weights gives the exact spectrum \[
\operatorname{Spec}D
=\left(
-\frac{2}{\sqrt{1+\epsilon^2}}^{[q-1]},
-\frac{1}{\sqrt{1+\epsilon^2}},0,
\frac{1}{\sqrt{1+\epsilon^2}},
\frac{2}{\sqrt{1+\epsilon^2}}^{[q-1]}
\right).
\tag{SR27}
\] Superscripts in square brackets denote multiplicities, as in SR2, rather than powers. The case \(q=1\) simply has one paired block of weight one and one zero coordinate; SR26 is \((c_3-c_1)/(1+\epsilon^2)\).

For any nonscalar prescribed spectrum, SR26 tends to \(B_q(c)\) as \(\epsilon\) decreases to zero through positive values. Yet \(\det M_\epsilon\) tends to zero by SR19, so the limiting form is outside the positive class. Together with SR13, this proves a sharp supremum and excludes an attained nontrivial maximum within the literal polynomial flags.

\subsubsection{7. Additional exact fixed-metric optimum for this family}\label{endpoint-four_window_sharpness_review--additional-exact-fixed-metric-optimum-for-this-family}

The source's chosen control in SR23 is independent of \(\epsilon\), and its trace is exactly SR26. The full difference spectrum SR27 also permits computing the optimum over all controls with the prescribed spectrum at any fixed \(\epsilon>0\).

Set \(\rho=\sqrt{1+\epsilon^2}\). On the \(j\)-th pair, define \[
f_j^+=\frac{(\rho+1)e_{j-1}+\epsilon e_{q+j}}
{\sqrt{2\rho(\rho+1)}},\qquad
f_j^-=\frac{-\epsilon e_{j-1}+(\rho+1)e_{q+j}}
{\sqrt{2\rho(\rho+1)}}.
\tag{SR28}
\] Their inner product is zero, and their squared norms are one because \((\rho+1)^2+\epsilon^2=2\rho(\rho+1)\). Multiplication by the unscaled matrix \(\begin{pmatrix}1&\epsilon\\\epsilon&-1\end{pmatrix}\) gives eigenvalues \(+\rho\) and \(-\rho\), respectively. Hence \(D_j f_j^\pm=\pm(w_j/\rho)f_j^\pm\).

Define \(\widetilde C_0\) to assign \(a_j\) to \(f_j^+\), \(b_j'\) to \(f_j^-\), and \(c_{q+1}\) to \(e_q\), and pull it back to the original polynomial space by \(T_\epsilon^{-1}\widetilde C_0T_\epsilon\). This is a fully specified \(M_\epsilon\)-self-adjoint operator with the required eigenvalue multiset. Its trace is \[
\operatorname{Tr}(D\widetilde C_\epsilon)
=\frac{1}{\rho}\sum_jw_j(a_j-b_j')
=\frac{B_q(c)}{\sqrt{1+\epsilon^2}}.
\tag{SR29}
\] For any other control with that spectrum, the spectral projection expansion SR11 and SR3 applied to \(D\) bound its trace above by the ordered pairing of its eigenvalues with SR27, which is the right side of SR29. Reversing the control assignment gives the corresponding negative bound. Thus the exact maximum of the absolute trace, at this fixed \(M_\epsilon\), is \(B_q(c)/\sqrt{1+\epsilon^2}\).

This additional formula preserves the correctness of the source's particular choice SR26. It computes the optimization left available when its prescribed spectrum is allowed a different self-adjoint realization. It remains in the same positive coefficient-metric family. Its full Ky Fan deficits are \(g_r(M_\epsilon)=a_r(1-1/\sqrt{1+\epsilon^2})\), by SR14 and SR27, where \(a_r\) here is the reference sum from SR9.

\subsubsection{8. Metric-path realization and its literal mass}\label{endpoint-four_window_sharpness_review--metric-path-realization-and-its-literal-mass}

The source's fixed polynomial-coordinate operator \(C_\epsilon\) satisfies \(C_\epsilon^*M_\epsilon=M_\epsilon C_\epsilon\). Therefore \[
M_\epsilon(t)
=e^{tC_\epsilon^*/2}M_\epsilon e^{tC_\epsilon/2}
=M_\epsilon e^{tC_\epsilon}
\tag{SR30}
\] is positive for every real \(t\). The equality follows by the exponential series and the stated intertwining. Differentiating gives \(M_\epsilon(t)^{-1}\dot M_\epsilon(t)=C_\epsilon\). Thus the operator is an actual logarithmic derivative, with the same eigenvalues, on an explicit positive coefficient-metric path.

For the locally affine path \(M_\epsilon+tM_\epsilon C_\epsilon\), positivity is equivalent to \(1+tc_i>0\) for all prescribed eigenvalues, after the same explicit isometry. Its logarithmic derivative at zero is again \(C_\epsilon\).

For clarity, the fixed mass statement \(\|1\|^2=\mu\) holds across the parameter family \(\epsilon>0\) at \(t=0\). The path's original constant polynomial has eigenvalue \(c_n\): this is \(a_1=c_{2q+1}\) for \(q\ge2\), and \(a_q=c_3\) at \(q=1\). Consequently its literal norm along SR30 is \(\mu e^{tc_n}\), and along the affine path is \(\mu(1+tc_n)\). No claim that those paths preserve mass in \(t\) is needed or made; these formulas give the actual factors.

\subsubsection{9. Exact moment obstruction and the retained arithmetic scope}\label{endpoint-four_window_sharpness_review--exact-moment-obstruction-and-the-retained-arithmetic-scope}

Put \(N=2q\). Let \(\mathcal M_N\) be the real vector space of finite real signed measures \(\nu\) satisfying \(\int(1+|u|^{2N})\,d|\nu|(u)<\infty\). On a vertical line \(S=k/2+iu\), the real-linear moment map \(\mathfrak M_k:\mathcal M_N\to\operatorname{Herm}(\mathcal P_N)\) is \[
\mathfrak M_k(\nu)(P,R)
=\int\overline{P(k/2+iu)}R(k/2+iu)\,d\nu(u).
\tag{SR31}
\] For a nonnegative measure the form is positive semidefinite, and it is positive definite when no nonzero polynomial in the source vanishes almost everywhere. The original positive densities satisfy this condition. For \(P,R\in\mathcal P_{N-1}\), evaluation gives the real-linear constraint map \(\mathcal L_k:\operatorname{Herm}(\mathcal P_N)\to\operatorname{Herm}(\mathcal P_{N-1})\), \[
\mathcal L_k(M)(P,R)
:=M(SP,R)+M(P,SR)-kM(P,R)=0
\quad\text{when }M=\mathfrak M_k(\nu).
\tag{SR32}
\] Indeed \(\overline{S(u)}+S(u)=k\) pointwise, so the three integrands cancel. In monomial coefficients, this linear map has entries \(M_{i+1,j}+M_{i,j+1}-kM_{ij}\) wherever both shifted indices occur in the finite domain. Thus SR31 has image in the kernel of the explicitly typed linear map \(\mathcal L_k\). This supplies the exact relation between moment forms and general coefficient forms.

In the constructed family take \(P=\chi,R=1\). All products belong to \(\mathcal P_{2q}\), because \(q+1\le2q\). Equations SR17--SR19 give \[
\langle\chi,1\rangle=0,\qquad
\langle S\chi,1\rangle=\mu.
\tag{SR33}
\] For \(q\ge2\), \(S=b_1\) is in the lower degree flag and is orthogonal to \(b_q=\chi\), so \(\langle\chi,S\rangle=0\). At \(q=1\), writing the original monic polynomial as \(\chi=S+\chi(0)\) gives \(S=\chi-\chi(0)\), hence \(\langle\chi,S\rangle=\mu\). The convention is conjugate-linear in the first slot and linear in the second; the coefficient \(-\chi(0)\) multiplies the zero cross pairing, so no unretained conjugation occurs. Therefore \[
\mathcal L_k(M_\epsilon)(\chi,1)
=\begin{cases}\mu,&q\ge2,\\2\mu,&q=1.\end{cases}
\tag{SR34}
\] This is nonzero for every \(\epsilon>0,\mu>0\), and the subtracted term is zero for every real center parameter \(k\). Thus this family violates a necessary moment identity for every real vertical line, not only the particular arithmetic center.

The original Gamma and arithmetic source forms remain in the image of their actual measures under SR31, with masses \((2\pi)^{k/2}\) and \(\mu_h^k\), and with the original \(\chi\), \(\eta\), and \(\sigma_h^{\otimes k}\eta\) from the accepted AT source. The constructed \(M_\epsilon\) is linked to its explicit coordinate model by SR17--SR19 and has the computed nonzero obstruction SR34. It therefore proves sharpness in the larger class of positive coefficient forms while proving no sharpness assertion for the original arithmetic moment class.

For each actual fixed positive moment Gram, the strict inequality SR13 still applies because its literal flags are the same and it belongs to the positive coefficient class. The approach-to-equality family supplied here cannot establish that the actual moment class approaches the same supremum: its obstruction SR34 persists at every finite parameter. No conclusion about a uniform arithmetic improvement or about the failure of the F1 or Tau Base program is drawn from this calibration.

\subsubsection{10. Final findings}\label{endpoint-four_window_sharpness_review--final-findings}

The strictness proof is valid for every nonscalar control, including arbitrary repeated spectra and \(q=1\). Both middle-rank equality cases and the complementary high-rank case are explicitly excluded by the two original polynomial intersections. The family preserves \(\chi\), its coefficient frame and relation flags, its base mass \(\mu\), every cross term, and all powers of \(\epsilon\). Its trace and full difference spectrum are correct. Its violation of the moment identity is exactly \(\mu\) or \(2\mu\) on the stated pair.

The supplementary independent calculation of FS20--FS35 agreed with these block formulas and required no correction. SR28--SR29 additionally compute the exact prescribed-spectrum optimum for each fixed metric in this family. All these conclusions are add-only; the previously sealed AT proof and artifacts remain unchanged.
